\chapter{Заморозка качественного родителя и границы массовой метрики}
\label{ch:v7-qualitative-parent-mass-metric-freeze-gate}

\section{Предмет заморозки}

Последовательность последних гейтов построила семейство положительных
действий на одном и том же 27-мерном вещественном срезе:
\begin{equation}
 \mathcal S_\eta=\mathcal S_E+\eta\|\mathcal R_U\|_{\rm HS}^2,
 \qquad \eta>0.
 \label{eq:v7-freeze-action-family}
\end{equation}
Здесь $\mathcal S_E$ есть рёберная Hodge-норма, а
$\mathcal R_U=C_tU-UC_s$ --- относительная кривизна канонического
полярного связывающего бимодуля.

Незавешенный представитель $\eta=1$ реализуется одним матричным следом
общей кривизны. Однако предыдущая глава доказала, что текущая порождённая
алгебра имеет нетривиальный центр и не фиксирует $\eta$. Поэтому предметом
заморозки является не одна количественно завершённая теория, а общий
качественный класс~\eqref{eq:v7-freeze-action-family}.

\section{Нормировочно независимое ядро}

В начале связывающая норма начинается с четвёртого порядка. Поэтому для
всех $\eta\geq0$
\begin{equation}
 \Hess_0\mathcal S_\eta=\Hess_0\mathcal S_E,
 \qquad
 (n_-,n_0,n_+)_0=(7,0,20).
 \label{eq:v7-freeze-origin-signature}
\end{equation}
Ровно семь корневых направлений запускаются, а двадцать down--weak
конкурентов сохраняют положительную щель $18/5$.

В целевой точке $A_0$ имеем
\begin{equation}
 \Hess_{A_0}\mathcal S_\eta
 =H_E+2\eta H_L,
 \qquad H_E>0,
 \qquad H_L\geq0,
 \label{eq:v7-freeze-vacuum-decomposition}
\end{equation}
причём $\rank H_L=22$. Следовательно,
\begin{equation}
 (n_-,n_0,n_+)_{A_0}=(0,0,27)
 \qquad\text{для всех }\eta\geq0.
 \label{eq:v7-freeze-vacuum-signature}
\end{equation}
Машинный контроль охватывает веса от $10^{-6}$ до $10^6$ и подтверждает
аналитический вывод. Таким образом, выбор семи запускающихся направлений и
поперечная устойчивость целевой точки не зависят от относительной метрики.

\begin{theorem}[Качественный универсальный класс Тома VII]
На исследованном вещественном срезе существует один общий
рангоизменяющий носитель и класс ограниченных снизу действий
$\mathcal S_\eta$, $\eta>0$. Каждый его представитель имеет стационарный
нулевой сектор с семью отрицательными и двадцатью положительными модами и
целевую стационарную точку со строго положительным 27-мерным гессианом.
Эти сигнатуры являются инвариантом класса и не требуют выбора массовой
нормировки.
\end{theorem}

\section{Почему это ещё не полный физический родитель}

Проверенный гессиан относится к семи корневым амплитудам и двадцати
тяжёлым конкурентам. Он не включает все ранее обнаруженные семейные и
фазовые степени свободы. После лесного факторирования сохраняется одна
голономия
\begin{equation}
 \mathcal M_C=U(3)/\operatorname{Ad}U(3),
 \label{eq:v7-freeze-phase-moduli}
\end{equation}
с тремя собственными фазами. Один классовый функционал этой голономии не
выбирает относительные семейные оси и потому не выводит CKM или PMNS.

Количественная метрика также остаётся свободной. Связывающая опора сама
порождает простой фактор $M_{21}(\mathbb C)$, но её объединение с рёберным
сектором является прямой суммой с центральным относительным весом.
Собственные значения $H_E+2\eta H_L$ меняются с $\eta$, хотя их знаки
сохраняются.

Наконец, одна размерная калибровка фиксирует только общий линейный масштаб.
Пространственно-временное произведение оставляет
\begin{equation}
 \lambda_E=\frac{\pi^2}{f_0},
 \label{eq:v7-freeze-f0-boundary}
\end{equation}
а единый физический калибровочный след, определяющий $f_0$, не построен.
Следовательно, абсолютные массы, нелинейные амплитуды и времена не являются
предсказаниями Тома VII.

\begin{nogocriterion}[Запрет расширительного чтения]
Нельзя переносить сигнатуры 27-мерного амплитудно-тяжёлого среза на полный
фазовый, калибровочный и пространственно-временной фактор без отдельного
гессиана. Нельзя выбирать $\eta$, $f_0$ или семейные фазы по наблюдаемым и
после этого называть полученные массы и смешивания выводом родителя.
Названия новых физических частиц до такого замыкания не вводятся.
\end{nogocriterion}

\section{Условия повторного открытия}

Количественная ветвь может быть открыта повторно только при появлении хотя
бы одного нового выведенного объекта:
\begin{enumerate}
 \item физического внедиагонального коннектора между рёберной и
 связывающей опорами;
 \item второго некоммутирующего семейного тензора, выбирающего
 относительные оси;
 \item одного общего калибровочно-пространственного следа, фиксирующего
 кинетическую и массовую метрики.
\end{enumerate}
Любой из этих объектов должен пройти Real-условие, первый порядок или его
явно зарегистрированную замену, полный физический гессиан и запрет
подгонки.

\begin{protocol}
\begin{enumerate}
 \item общий исправленный носитель: \textbf{получен};
 \item один незавешенный представитель действия: \textbf{получен};
 \item качественный класс $\mathcal S_\eta$: \textbf{получен};
 \item нулевая сигнатура: \textbf{$(7,0,20)$};
 \item вакуумная сигнатура на проверенном срезе: \textbf{$(0,0,27)$};
 \item устойчивость сигнатур при всех $\eta>0$: \textbf{доказана};
 \item единственная массовая метрика: \textbf{не выведена};
 \item полный фазовый гессиан: \textbf{не получен};
 \item CKM/PMNS: \textbf{не выведены};
 \item абсолютный масштаб и квартика: \textbf{не замкнуты};
 \item полный физический родитель: \textbf{не получен};
 \item исследовательский вопрос Тома VII: \textbf{завершён};
 \item численные массовые предсказания: \textbf{запрещены};
 \item итог: \textbf{качественный класс закрыт, физическое замыкание нет}.
\end{enumerate}
\end{protocol}

Машинный реестр сохранён в
\path{s2t_v7_qualitative_parent_mass_metric_freeze_gate_results.json}.